\uuid{Kv8O}
\exo7id{6295}
\titre{exo7 6295}
\auteur{queffelec}
\organisation{exo7}
\datecreate{2011-10-16}
\isIndication{false}
\isCorrection{false}
\chapitre{Autre}
\sousChapitre{Autre}
\cpge{bcpst1.derivees-primitives-integrales.b, bcpst1.fonctions-deux-variables.c, mp-mpi.calcul-differentiel-optimisation.a, mp-mpi.calcul-differentiel-optimisation.b, pc.fonctions-plusieurs-variables.c, psi.fonctions-plusieurs-variables.c}
\module{Calcul différentiel}
\niveau{L3}
\difficulte{}

\contenu{
\texte{
Soit $U$ l'ouvert $\Rr^3$ privé de l'axe des
$z$.
}
\begin{enumerate}
    \item \question{Vérifiez que la fonction $f(x,y,z)$, qui vaut $e^z
\cos {\theta\over 2}{\sin r\over \sqrt r}$ en
coordonnées cylindriques, est harmonique sur $U$.}
    \item \question{Soit $\lambda $ une constante réelle. Montrer qu'une
fonction du type $\displaystyle f(x,y,z) = e^z \cos
(\lambda \theta) u(r)$ est harmonique dans $U$ si et
seulement si $u = u(r)$ est solution de l'équation
différen\-tielle (dite de \emph{Bessel}) :
\begin{equation*}r^2 u''(r) + r u'(r) + [r^2 - 
\lambda^2] u (r) = 0.\tag*{$(E_\lambda)$}\end{equation*}}
    \item \question{Vérifiez, que pour $\lambda = 3/2$, la fonction
$\displaystyle u(r) = {\sin r - r \cos r\over r \sqrt r}$
convient.}
\end{enumerate}
}
